\documentclass[11pt]{article}
\usepackage[a4paper,margin=1in]{geometry}
\usepackage{fontspec}
\usepackage[boldfont=false]{xeCJK}
\setCJKmainfont{FandolSong-Regular.otf}
\usepackage{amsmath}
\usepackage{amssymb}
\usepackage{array}
\usepackage{booktabs}
\usepackage{graphicx}
\usepackage{adjustbox}
\usepackage{enumitem}
\setlist[itemize]{topsep=2pt,partopsep=0pt,parsep=0pt,itemsep=2pt,leftmargin=1.4em}
\setlist[enumerate]{topsep=2pt,partopsep=0pt,parsep=0pt,itemsep=2pt,leftmargin=1.6em}
\usepackage{parskip}
\usepackage{fvextra}
\fvset{breaklines=true,breakanywhere=true,fontsize=\small}
\setlength{\parindent}{0pt}

\begin{document}

\begin{center}
  {\LARGE\bfseries Inheritance}\\[0.35em]
  {\large Igcse Biology}
\end{center}
\vspace{0.6em}

\section*{Chromosomes, genes and DNA}

\textbf{Chromosomes} 染色体 are made of \textbf{DNA}, which carries the genetic information in units called \textbf{genes} 基因.

\begin{itemize}
\item A \textbf{gene} is a length of DNA that \textbf{codes} 编码 for one \textbf{protein} 蛋白质.

\item An \textbf{allele} 等位基因 is an alternative form of a gene. For example, a gene for eye colour may have a brown allele and a blue allele.

\end{itemize}

\subsection*{Sex chromosomes}

Sex is inherited through the \textbf{sex chromosomes} 性染色体. Females are \textbf{XX} and males are \textbf{XY}. An egg always carries an X; a sperm carries either an X or a Y, so the sperm decides the baby's sex. This gives a 1:1 \textbf{ratio} 比例 of girls to boys.

\section*{How DNA controls the cell (Supplement)}

The \textbf{base sequence} 碱基序列 of a gene sets the order of \textbf{amino acids} 氨基酸 in a protein. A different order of amino acids folds the protein into a different shape, and the shape decides its job. So DNA controls the cell by controlling which proteins are made — including \textbf{enzymes} 酶, membrane carriers and \textbf{receptor proteins} 受体蛋白 for \textbf{neurotransmitters} 神经递质.

Most body cells hold the \textbf{same} genes, but each cell only \textbf{expresses} 表达 (switches on) the genes it needs, so it makes only the proteins for its own job.

\section*{Making a protein (Supplement)}

A gene stays in the \textbf{nucleus} 细胞核, but proteins are made in the \textbf{cytoplasm} 细胞质. The link between them is \textbf{messenger} 信使 RNA (mRNA):

\begin{enumerate}
\item mRNA is made in the nucleus as a copy of the gene.

\item the mRNA moves out into the cytoplasm and passes through a \textbf{ribosome} 核糖体.

\item the ribosome reads the mRNA's \textbf{bases} 碱基 and joins amino acids in the matching order to build the protein.

\end{enumerate}

\section*{Cell division (Supplement)}

Body cells are \textbf{diploid} 二倍体 (two sets of chromosomes; humans have 23 pairs). Gametes are \textbf{haploid} 单倍体 (one set). Cells make new cells by nuclear \textbf{division} 分裂, of two kinds.

\textbf{Mitosis} 有丝分裂:

\begin{itemize}
\item makes two cells that are genetically \textbf{identical} to the parent cell.

\item the chromosomes are copied exactly (\textbf{replication} 复制) before division, so each \textbf{daughter cell} 子细胞 keeps the full chromosome number.

\item it is used for growth, repair of tissues, replacing old cells, and \textbf{asexual reproduction} 无性生殖.

\item \textbf{stem cells} 干细胞 are unspecialised cells that divide by mitosis; their daughter cells can then become specialised.

\end{itemize}

\textbf{Meiosis} 减数分裂:

\begin{itemize}
\item makes \textbf{gametes} 配子.

\item it is a \textbf{reduction division}: the chromosome number is halved, from diploid to haploid.

\item the cells it makes are genetically \textbf{different} from one another.

\end{itemize}

\section*{Inheritance: key words}

\textbf{Inheritance} 遗传 is the passing of genetic information from one generation to the next. Learn these words:

\begin{center}
\adjustbox{max width=\linewidth}{%
\begin{tabular}{ll}
\toprule
\textbf{Word} & \textbf{Meaning} \\
\midrule
\textbf{genotype} 基因型 & the alleles an organism has (its genetic make-up) \\
\textbf{phenotype} 表现型 & the features you can observe \\
\textbf{dominant} 显性 allele & shown in the phenotype even if only one copy is present (written as a capital letter, e.g. B) \\
\textbf{recessive} 隐性 allele & shown only when \textbf{no} dominant allele is present (small letter, e.g. b) \\
\textbf{homozygous} 纯合子 & two \textbf{identical} alleles (BB or bb); these breed true, or \textbf{pure-breeding} 纯种 \\
\textbf{heterozygous} 杂合子 & two \textbf{different} alleles (Bb); not pure-breeding \\
\bottomrule
\end{tabular}}
\end{center}

\section*{Monohybrid crosses}

To predict the offspring of a \textbf{cross} 杂交 involving one gene, use a genetic diagram or a \textbf{Punnett square} 庞纳特方格:

\begin{itemize}
\item write each parent's genotype, then the gametes (each gamete carries one allele).

\item combine them in a grid to find the possible offspring genotypes and phenotypes.

\end{itemize}

Two heterozygous parents (Bb × Bb) give a \textbf{3:1} ratio of dominant to recessive. A heterozygous crossed with a homozygous recessive (Bb × bb) gives a \textbf{1:1} ratio. You can also read a \textbf{pedigree diagram} 系谱图 (a family tree) to follow a feature and work out genotypes.

\textbf{(Supplement)} A \textbf{test cross} 测交 finds an unknown genotype: cross the unknown with a homozygous recessive. If any offspring show the recessive feature, the unknown must have been heterozygous.

\section*{Codominance and blood groups (Supplement)}

In \textbf{codominance} 共显性, \textbf{both} alleles in a heterozygous organism show in the phenotype.

The ABO \textbf{blood groups} 血型 are an example. The alleles are $I^A$, $I^B$ and $I^o$. The alleles $I^A$ and $I^B$ are codominant, while $I^o$ is recessive:

\begin{itemize}
\item $I^A I^A$ or $I^A I^o$ → group A

\item $I^B I^B$ or $I^B I^o$ → group B

\item $I^A I^B$ → group AB (both alleles shown)

\item $I^o I^o$ → group O

\end{itemize}

\section*{Sex linkage (Supplement)}

A \textbf{sex-linked} 伴性遗传 characteristic is controlled by a gene on a sex chromosome (usually the X). Because males have only one X chromosome, a recessive allele on it always shows in males, so the feature is more common in males than in females.

\textbf{Red-green colour blindness} 色盲 is an example: the recessive allele is on the X chromosome, so it is much more common in boys than in girls.

\section*{Exam tips}

\begin{itemize}
\item Gene = a length of DNA coding for a protein. Allele = one form of a gene.

\item Genotype = the alleles present; phenotype = what you see. Homozygous = two same; heterozygous = two different.

\item Dominant shows with one copy; recessive needs two copies.

\item Learn to draw a Punnett square. Bb × Bb → 3:1; Bb × bb → 1:1.

\item (Supplement) Codominance: both alleles show (blood group AB). Sex linkage: a gene on the X, more common in boys (colour blindness).

\end{itemize}


\end{document}
